%!TEX root = ../../main.tex
%% ===============================================================
%% 2.4 평가 지표의 정의와 기원
%% 파일: chapters/02_background/04_metrics.tex
%% 상위: chapters/02_background/chapter.tex
%% 정의 라벨: sec:bg-metrics, eq:auc
%% 참조 라벨: eq:auc
%% 포함 객체: -
%% 인용: brier1950verification, davis2006relationship, fawcett2006introduction, green1966signal, guo2017calibration, hanley1982meaning, mann1947test, naeini2015obtaining
%% 상태: 초안 (docs/OUTLINE.md 참고)
%% ===============================================================

\section{평가 지표의 정의와 기원}
\label{sec:bg-metrics}

\begin{definition}[ROC 곡선과 AUC]
임계값 $\tau$에 대해 참양성률 $\mathrm{TPR}(\tau)=P(f(X)\ge\tau \mid Y=1)$과 거짓양성률 $\mathrm{FPR}(\tau)=P(f(X)\ge\tau \mid Y=0)$을 정의하면, ROC(receiver operating characteristic) 곡선은 $\tau$를 변화시키며 얻는 점 $(\mathrm{FPR}(\tau), \mathrm{TPR}(\tau))$의 궤적이다. AUC는 그 아래 면적으로,
\begin{equation}
  \AUC = P\big(f(X^{+}) > f(X^{-})\big) + \tfrac12 P\big(f(X^{+}) = f(X^{-})\big)
  \label{eq:auc}
\end{equation}
와 같이 무작위로 뽑은 양성 표본의 점수가 음성 표본의 점수보다 클 확률과 같다.
\end{definition}

ROC는 제2차 세계대전기의 레이더 신호 검출 이론에서 유래하였고 심리물리학으로 확장되었다 \cite{green1966signal}. 식 \eqref{eq:auc}의 확률 해석은 Hanley와 McNeil \cite{hanley1982meaning}이 AUC가 Mann--Whitney $U$ 통계량 \cite{mann1947test}과 동치임을 보이며 정립하였다. 기계학습에서의 표준적 소개는 Fawcett \cite{fawcett2006introduction}을 따른다. 본 논문은 교전 라벨이 거의 균형($\approx 50/50$)이고, 예측기의 순위 품질이 관심사이므로 AUC를 1차 지표로 사용한다.

\begin{definition}[평균 정밀도]
정밀도 $\mathrm{Prec}(\tau)$와 재현율 $\mathrm{Rec}(\tau)$의 곡선 아래 면적으로, 임계값 순서대로 $\mathrm{AP} = \sum_k \big(\mathrm{Rec}_k - \mathrm{Rec}_{k-1}\big)\mathrm{Prec}_k$로 계산한다. 정밀도--재현율 곡선은 불균형 상황에서 ROC보다 정보량이 크다 \cite{davis2006relationship}.
\end{definition}

\begin{definition}[Brier 점수]
$N$개의 예측 확률 $p_i$와 결과 $y_i$에 대해 $\mathrm{BS} = \frac{1}{N}\sum_{i=1}^{N}(p_i - y_i)^2$. 1950년 Brier가 기상 예보의 검증을 위해 제안하였으며 \cite{brier1950verification}, 엄격한 적정 채점 규칙(strictly proper scoring rule)이다.
\end{definition}

\begin{definition}[보정과 기대 보정 오차]
예측 확률이 \emph{보정}(calibrated)되었다는 것은 $P(Y=1 \mid p(X)=p) = p$가 모든 $p$에 대해 성립함을 뜻한다. 예측을 $M$개의 구간 $B_m$으로 나누었을 때 기대 보정 오차는
\begin{equation}
  \mathrm{ECE} = \sum_{m=1}^{M} \frac{|B_m|}{N}\,\Big|\,\mathrm{acc}(B_m) - \mathrm{conf}(B_m)\Big|
\end{equation}
이다. 구간 기반 ECE는 Naeini 등 \cite{naeini2015obtaining}이 제안하였고, 심층 신경망의 과신(overconfidence)과 온도 조정(temperature scaling)에 의한 사후 보정은 Guo 등 \cite{guo2017calibration}이 체계화하였다.
\end{definition}
